% Aula 09 --- o ou exclusivo: convexidade, a ampliação pelos quadrados e o termo cruzado.
Mariana estuda engenharia e quer que um classificador linear preveja quando uma lâmpada
acende. A lâmpada tem duas chaves, $x_1,x_2\in\{0,1\}$, e acende ($y=+1$) quando exatamente
uma delas está ligada; com as duas desligadas ou as duas ligadas, fica apagada ($y=-1$). As
quatro observações são $(0,0)$ e $(1,1)$, na classe $-1$, e $(1,0)$ e $(0,1)$, na classe
$+1$. Lembre que observações $(\bm z_i,y_i)$ de $\R^m$ são linearmente separáveis se existe
$(\beta_0,\bbeta)$, com $\bbeta\ne\bm0$, tal que $y_i(\beta_0+\bbeta^\top\bm z_i)>0$ para todo
$i$, e que, nesse caso, o classificador de margem máxima é dado pela solução de
\[
  \max_{\beta_0,\,\bbeta,\,M}\ M
  \quad\text{sujeito a}\quad
  \norm{\bbeta}=1
  \quad\text{e}\quad
  y_i\,(\beta_0+\bbeta^\top\bm z_i)\ge M\ \text{ para todo } i .
\]
\begin{itens}
  \item \pts{1,0} Suponha que algum ponto $\bm z\in\R^m$ seja combinação convexa de
        observações da classe $+1$ e também combinação convexa de observações da classe
        $-1$. Mostre que as observações não são linearmente separáveis. Por que isso vale,
        em particular, quando uma observação de uma classe está no interior de um triângulo
        formado por observações da outra?
  \item \pts{0,6} Mostre que as quatro observações de Mariana não são linearmente separáveis.
  \item \pts{1,0} Mariana segue a primeira ideia das notas para obter fronteiras não lineares:
        acrescenta os quadrados e trabalha com as covariáveis $(x_1,x_2,x_1^2,x_2^2)$. Mostre
        que as observações continuam não separáveis nesse espaço ampliado.
  \item \pts{1,4} Ela troca os quadrados pelo produto e usa as três covariáveis
        $\bm z=(x_1,\,x_2,\,x_1x_2)$. Mostre que $f(\bm z)=2z_1+2z_2-4z_3-1$ satisfaz
        $y_if(\bm z_i)=1$ nas quatro observações. Mostre que o classificador de margem
        máxima nesse espaço é o hiperplano $f(\bm z)=0$, com $M=1/\sqrt{24}\approx0{,}204$,
        e que as quatro observações são vetores de suporte. (Para a cota superior, multiplique
        as restrições das observações $(0,0)$, $(1,1)$, $(1,0)$ e $(0,1)$ pelos pesos $8$,
        $4$, $6$ e $6$, some e use a desigualdade de Cauchy--Schwarz.) Que fronteira esse
        hiperplano define no plano $(x_1,x_2)$?
\end{itens}

\begin{solucao}
\textbf{(a)} Escreva $\bm z=\sum_j\lambda_j\bm z_j^+=\sum_k\mu_k\bm z_k^-$, com
$\bm z_j^+$ observações da classe $+1$, $\bm z_k^-$ observações da classe $-1$, pesos
$\lambda_j,\mu_k\ge0$ e $\sum_j\lambda_j=\sum_k\mu_k=1$. Suponha, por absurdo, que
$f(\bm z)=\beta_0+\bbeta^\top\bm z$ separe as classes: $f(\bm z_j^+)>0$ para todo $j$ e
$f(\bm z_k^-)<0$ para todo $k$. Como os pesos somam $1$,
\[
  f\Big(\sum_j\lambda_j\bm z_j^+\Big)=\sum_j\lambda_j\beta_0+\sum_j\lambda_j\bbeta^\top\bm z_j^+
  =\sum_j\lambda_jf(\bm z_j^+) ,
\]
e uma combinação convexa de números positivos é positiva: $f(\bm z)>0$. Pela mesma conta com
os $\mu_k$, $f(\bm z)<0$, o que é absurdo. Em termos geométricos, cada lado de um hiperplano é
convexo, e um ponto que é combinação convexa das duas classes teria de estar dos dois lados
ao mesmo tempo. Se uma observação $\bm z^+$ está no interior de um triângulo de observações
da classe $-1$, ela é combinação convexa dos três vértices, e também, trivialmente, de si
mesma, com peso $1$: a hipótese vale com $\bm z=\bm z^+$.

\smallskip
\textbf{(b)} O centro do quadrado é combinação convexa das duas classes:
\[
  \big(\tfrac12,\tfrac12\big)=\tfrac12(1,0)+\tfrac12(0,1)=\tfrac12(0,0)+\tfrac12(1,1) .
\]
Pelo item (a), nenhuma reta separa as observações.

\smallskip
\textbf{(c)} Em $\{0,1\}$, $x_j^2=x_j$. Cada observação ampliada é, então,
$T(\x)=(x_1,x_2,x_1,x_2)$, em que $T:\R^2\to\R^4$ é linear. Uma aplicação linear preserva
combinações convexas, e, aplicando $T$ à identidade do item (b),
\[
  T\big(\tfrac12,\tfrac12\big)=\tfrac12T(1,0)+\tfrac12T(0,1)=\tfrac12T(0,0)+\tfrac12T(1,1),
\]
um ponto de $\R^4$ que é combinação convexa das duas classes ampliadas. Pelo item (a), que vale
em qualquer dimensão, as observações ampliadas não são separáveis. Dito de outro modo, um
hiperplano de $\R^4$ calculado nas observações ampliadas vale
$\beta_0+(\beta_1+\beta_3)x_1+(\beta_2+\beta_4)x_2$: é uma função afim de $(x_1,x_2)$, que o
item (b) já mostrou incapaz de separar as classes. Os quadrados não acrescentam nada a
variáveis binárias.

\smallskip
\textbf{(d)} As observações ampliadas são $(0,0,0)$ e $(1,1,1)$, na classe $-1$, e
$(1,0,0)$ e $(0,1,0)$, na classe $+1$, e
\[
  f(0,0,0)=-1,\qquad f(1,1,1)=2+2-4-1=-1,\qquad f(1,0,0)=f(0,1,0)=2-1=1 :
\]
$y_if(\bm z_i)=1$ nas quatro. Dividindo os coeficientes por $\norm{(2,2,-4)}=\sqrt{24}$,
obtém-se $\bbeta=(2,2,-4)/\sqrt{24}$, de norma $1$, e $\beta_0=-1/\sqrt{24}$, com
$y_i(\beta_0+\bbeta^\top\bm z_i)=1/\sqrt{24}$ para todo $i$: o valor $M=1/\sqrt{24}$ é
atingido.

Para a cota, seja $(\beta_0,\bbeta,M)$ qualquer ponto que satisfaça as restrições.
Multiplicando a restrição de cada observação pelo seu peso $c_i>0$ e somando,
\[
  \beta_0\sum_ic_iy_i+\bbeta^\top\Big(\sum_ic_iy_i\bm z_i\Big)\ \ge\ M\sum_ic_i .
\]
Com os pesos $8$, $4$, $6$ e $6$, $\sum_ic_iy_i=-8-4+6+6=0$, e $\beta_0$ desaparece;
$\sum_ic_iy_i\bm z_i=-4(1,1,1)+6(1,0,0)+6(0,1,0)=(2,2,-4)$; e $\sum_ic_i=24$. Pela
desigualdade de Cauchy--Schwarz e por $\norm{\bbeta}=1$,
\[
  24\,M\ \le\ \bbeta^\top(2,2,-4)\ \le\ \norm{(2,2,-4)}=\sqrt{24},
  \qquad\text{isto é,}\qquad M\le\frac1{\sqrt{24}} .
\]
O ponto encontrado acima é, portanto, ótimo. É também a única solução: a igualdade em
Cauchy--Schwarz exige $\bbeta=(2,2,-4)/\sqrt{24}$, e, como todos os pesos são positivos, a
igualdade na soma exige que as quatro restrições valham com igualdade, o que determina
$\beta_0=-1/\sqrt{24}$. As quatro observações estão sobre a margem: todas são vetores de
suporte.

Os pesos não são arbitrários. Como $\sum_ic_iy_i\bm z_i=(2,2,-4)$, vale
$f(\bm z)=-1+\sum_ic_iy_i\inner{\bm z}{\bm z_i}$: os números $c_iy_i$, isto é, $-8$, $-4$, $6$
e $6$, dão uma representação de $f$ por produtos internos com as observações, como a que as
notas admitem sem demonstração. Ela não é a única, porque $\inner{\bm z}{\bm 0}=0$ e o
coeficiente da observação $(0,0,0)$ pode ser qualquer; o valor $-8$ é o que faz os
coeficientes somarem zero, e é essa soma nula que elimina $\beta_0$ da cota.

No plano original, $f=2x_1+2x_2-4x_1x_2-1=-4\big(x_1-\tfrac12\big)\big(x_2-\tfrac12\big)$. A
fronteira $f=0$ é o par de retas $x_1=\tfrac12$ e $x_2=\tfrac12$, e o classificador acende a
lâmpada quando $(x_1-\tfrac12)(x_2-\tfrac12)<0$, isto é, quando exatamente uma das chaves está
acima de $\tfrac12$. O que o ou exclusivo pede é a interação $x_1x_2$, e não os quadrados. O
\emph{kernel} polinomial de grau $2$, $\big(1+\inner{\x}{\x'}\big)^2$, contém o monômio $x_1x_2$
entre as coordenadas da sua função $\varphi$, e uma SVM com ele também separa estas
observações, embora com coeficientes diferentes, porque as coordenadas de $\varphi$ vêm
multiplicadas por constantes.
\end{solucao}
